\chapter{Triển khai}

\section{Mô hình triển khai}

\begin{figure}[H]
\centering
\begin{tikzpicture}[
    node distance=4cm and 3cm,
    every node/.style={font=\small},
    controller/.style={
        rectangle, rounded corners,
        minimum width=3cm, minimum height=1.5cm,
        text centered, draw=black, fill=red!20,
        align=center
    },
    worker/.style={
        rectangle, rounded corners,
        minimum width=2.8cm, minimum height=1.5cm,
        text centered, draw=black, fill=blue!20,
        align=center
    },
    arrow/.style={->, very thick}
]

\node (worker1) [worker] {Worker 1\\Ubuntu 22.04};
\node (worker2) [worker, below=1cm of worker1] {Worker 2\\Ubuntu 24.04};
\node (controller) [controller, left=4cm of $(worker1)!0.5!(worker2)$] {Controller\\Arch Linux};

\draw [arrow] (controller.east) -- (worker1.west);
\draw [arrow] (controller.east) -- (worker2.west);

\end{tikzpicture}
\caption{Mô hình triển khai Controller--Workers}
\label{fig:deployment_model}
\end{figure}

\noindent Hệ thống được triển khai theo mô hình \textit{Controller--Workers}, trong đó máy Controller đóng vai trò trung tâm điều phối, thực thi các playbook Ansible và quản lý cấu hình. Các Worker là các máy mục tiêu nhận cấu hình thông qua kết nối SSH và thực thi các task tương ứng.

\noindent Mô hình này cho phép quản lý tập trung, hỗ trợ triển khai đồng thời trên nhiều hệ thống và đảm bảo tính nhất quán trong quá trình cấu hình.

\begin{table}[H]
\centering
\small
\resizebox{1.0\textwidth}{!}{
\begin{tabular}{|p{2cm}|p{2cm}|p{2.5cm}|p{7cm}|}
\hline
\textbf{Tên} & \textbf{Phiên bản} & \textbf{IP} & \textbf{Các gói cài đặt} \\
\hline
Controller & 2026.02.01 & 192.168.1.10 & 
\textbf{Core:} Ansible \\
& & & \textbf{KVM:} qemu-kvm, libvirt-daemon-system, libvirt-clients, bridge-utils, virtinst \\
& & & \textbf{Cloud-init:} cloud-image-utils, genisoimage \\
& & & \textbf{Base images:} ubuntu-22.04.qcow2, ubuntu-24.04.qcow2 \\
\hline
Worker 1 & Ubuntu 22.04 & 192.168.1.11 & 
\textbf{Ban đầu:} OpenSSH, Python3 \\
& & & \textbf{Phase 1:} apparmor, apparmor-utils \\
& & & \textbf{Phase 2:} fail2ban, libpam-tmpdir, debsums \\
& & & \textbf{Phase 5:} libpam-modules, libpam-pwquality, sudo \\
& & & \textbf{Phase 6:} auditd, audispd-plugins, rsyslog, aide \\
& & & \textbf{Phase 7:} systemd-journal-remote \\
& & & \textbf{Security:} lynis, chkrootkit, rkhunter, wget, unzip \\
\hline
Worker 2 & Ubuntu 24.04 & 192.168.1.12 & 
\textbf{Ban đầu:} OpenSSH, Python3 \\
& & & \textbf{Phase 1:} apparmor, apparmor-utils \\
& & & \textbf{Phase 2:} fail2ban, libpam-tmpdir, debsums \\
& & & \textbf{Phase 5:} libpam-modules, libpam-pwquality, sudo \\
& & & \textbf{Phase 6:} auditd, audispd-plugins, rsyslog, aide \\
& & & \textbf{Phase 7:} systemd-journal-remote \\
& & & \textbf{Security:} lynis, chkrootkit, rkhunter, wget, unzip \\
\hline
\end{tabular}
}
\caption{Thông tin chi tiết các máy trong mô hình và các gói được cài đặt}
\label{tab:system_info}
\end{table}

% =========================

\subsection{Quy trình Automated Security Hardening}

\noindent Dựa trên mô hình triển khai đã trình bày ở Hình~\ref{fig:deployment_model}, quy trình \textit{Automated Security Hardening} được xây dựng nhằm tự động hóa việc triển khai và đánh giá bảo mật hệ thống.

\noindent Quy trình hoạt động của hệ thống diễn ra theo chu trình khép kín:

\begin{itemize}
\item Auto-VM khởi tạo môi trường và cung cấp các Worker sẵn sàng
\item CIS-Ubuntu thực hiện hardening theo tiêu chuẩn bảo mật
\item Report Module đánh giá và phát hiện các điểm chưa đạt
\item Kết quả được phản hồi để cải tiến cấu hình trong các lần triển khai tiếp theo
\end{itemize}

\noindent Cách tiếp cận này giúp tự động hóa toàn bộ vòng đời bảo mật hệ thống, từ khởi tạo, cấu hình đến đánh giá và tối ưu hóa.

\begin{figure}[H]
\centering
\begin{tikzpicture}[
    node distance=5cm and 3cm,
    every node/.style={font=\small},
    component/.style={
        rectangle, rounded corners,
        minimum width=3.5cm, minimum height=2cm,
        text centered, draw=black,
        align=center, line width=0.8pt
    },
    auto_vm/.style={component, fill=green!20},
    cis_ubuntu/.style={component, fill=blue!20},
    report/.style={component, fill=orange!20},
    arrow/.style={->, very thick, line width=0.8pt},
    feedback/.style={->, very thick, dashed, bend right=45, line width=0.8pt}
]

\node (auto_vm) [auto_vm] {Auto-VM\\Tạo và cấu hình\\máy ảo};
\node (cis_ubuntu) [cis_ubuntu, right=of auto_vm] {CIS-Ubuntu\\Hardening\\bảo mật};
\node (report) [report, right=of cis_ubuntu] {Đánh giá và\\báo cáo};

\draw [arrow] (auto_vm.east) -- node[above] {VM sẵn sàng} (cis_ubuntu.west);
\draw [arrow] (cis_ubuntu.east) -- node[above] {Hardened} (report.west);

\draw [feedback] (report.north) to node[above] {Phản hồi cải tiến} (auto_vm.north);

\node[below=0.3cm of auto_vm, font=\tiny] {KVM/QEMU, Ansible, Cloud-init};
\node[below=0.3cm of cis_ubuntu, font=\tiny] {CIS Ubuntu 24.04 Benchmark};
\node[below=0.3cm of report, font=\tiny] {Lynis, OpenSCAP};

\end{tikzpicture}
\caption{Mô hình Automated Security Hardening System}
\label{fig:automated_security_model}
\end{figure}

\section{Manual}

\subsection{Initial Setup}

\noindent \textbf{Cài đặt các gói bảo mật cơ bản:}
\begin{verbatim}
sudo apt update && sudo apt upgrade -y                   # Cập nhật hệ thống
sudo apt install -y apparmor apparmor-utils \            # Framework bảo mật AppArmor
                   fail2ban \                            # Bảo vệ SSH khỏi brute-force
                   debsums \                             # Kiểm tra integrity packages
                   auditd audispd-plugins \              # Hệ thống audit logs
                   aide \                                # Intrusion detection system
                   rsyslog \                             # Centralized logging
                   libpam-pwquality \                    # Password quality enforcement
                   libpam-modules \                      # PAM authentication modules
                   sudo \                                # Privilege management
                   systemd-journal-remote                # Remote journal logging
\end{verbatim}

\noindent \textbf{Cấu hình AppArmor:}
\begin{verbatim}
sudo systemctl enable apparmor                           # Kích hoạt AppArmor boot
sudo systemctl start apparmor                            # Khởi động AppArmor service
sudo aa-enforce /etc/apparmor.d/*                        # Enforce tất cả profiles

sudo tee -a /etc/default/grub > /dev/null <<EOF
apparmor=1 security=apparmor
EOF

sudo update-grub                                         # Cập nhật bootloader
\end{verbatim}

\noindent \textbf{Cấu hình filesystem security:}
\begin{verbatim}
sudo mount -o remount,nodev,nosuid,noexec /tmp           # Bảo vệ /tmp directory
sudo mount -o remount,nodev,nosuid,noexec /dev/shm       # Bảo vệ shared memory

sudo tee -a /etc/fstab > /dev/null <<EOF
tmpfs /var/tmp tmpfs defaults,nodev,nosuid,noexec 0 0
EOF

sudo mount -a                                            # Mount lại filesystems
\end{verbatim}

\noindent \textbf{Cấu hình login banners:}
\begin{verbatim}
sudo tee /etc/issue /etc/issue.net /etc/motd > /dev/null <<EOF
Authorized uses only. All activity may be monitored and reported.
EOF

sudo chmod 644 /etc/issue /etc/issue.net /etc/motd       # Set permissions

sudo tee -a /etc/ssh/sshd_config > /dev/null <<EOF
Banner /etc/issue.net
EOF

sudo systemctl restart sshd                              # Restart SSH service
\end{verbatim}

\noindent \textbf{Kết quả mong đợi:}
\begin{itemize}
\item Hệ thống được cập nhật với các bản vá bảo mật mới nhất
\item AppArmor được kích hoạt và enforce tất cả profiles
\item Các thư mục tạm thời (/tmp, /dev/shm, /var/tmp) được bảo vệ với các tùy chọn bảo mật
\item Login banners hiển thị cảnh báo bảo mật khi người dùng đăng nhập
\item Các gói bảo mật cơ bản được cài đặt và sẵn sàng cấu hình
\end{itemize}

\subsection{Services Configuration}

\noindent \textbf{Disable unnecessary services:}
\begin{verbatim}
sudo systemctl stop telnet.socket telnet \               # Disable telnet services
sudo systemctl disable telnet.socket telnet
sudo systemctl stop rsh.socket rsh \                     # Disable rsh services
sudo systemctl disable rsh.socket rsh
sudo systemctl stop rlogin.socket rlogin \               # Disable rlogin services
sudo systemctl disable rlogin.socket rlogin
sudo systemctl stop ypbind && \                          # Disable NIS service
sudo systemctl disable ypbind
sudo systemctl stop tftp.socket tftp \                   # Disable tftp services
sudo systemctl disable tftp.socket tftp
sudo systemctl stop isc-dhcp-server && \                 # Disable DHCP server
sudo systemctl disable isc-dhcp-server
sudo systemctl stop slapd && \                           # Disable OpenLDAP
sudo systemctl disable slapd
sudo systemctl stop apache2 && \                         # Disable Apache web server
sudo systemctl disable apache2
sudo systemctl stop nfs-server && \                      # Disable NFS server
sudo systemctl disable nfs-server
sudo systemctl stop rpcbind && \                         # Disable RPC service
sudo systemctl disable rpcbind
sudo systemctl stop cups && \                            # Disable printing service
sudo systemctl disable cups
sudo systemctl stop avahi-daemon && \                    # Disable mDNS responder
sudo systemctl disable avahi-daemon
sudo systemctl stop bluetooth && \                       # Disable Bluetooth service
sudo systemctl disable bluetooth
\end{verbatim}

\noindent \textbf{Configure fail2ban for intrusion prevention:}
\begin{verbatim}
sudo apt install -y fail2ban                             # Install fail2ban
sudo systemctl enable fail2ban                           # Enable fail2ban boot
sudo systemctl start fail2ban                            # Start fail2ban service

cat << EOF | sudo tee /etc/fail2ban/jail.local           # Configure fail2ban
[DEFAULT]
bantime = 3600
findtime = 600
maxretry = 3

[sshd]
enabled = true
port = ssh
logpath = /var/log/auth.log
maxretry = 3
EOF

sudo systemctl restart fail2ban                          # Apply configuration
\end{verbatim}

\noindent \textbf{Configure time synchronization:}
\begin{verbatim}
sudo apt install -y ntp                                  # Install NTP service
sudo systemctl enable ntp                                # Enable NTP boot
sudo systemctl start ntp                                 # Start NTP service
sudo timedatectl set-ntp true                            # Enable system time sync
\end{verbatim}

\noindent \textbf{Kết quả mong đợi:}
\begin{itemize}
\item Các dịch vụ không cần thiết (telnet, rsh, rlogin, etc.) được vô hiệu hóa hoàn toàn
\item Fail2ban được cấu hình để bảo vệ SSH khỏi tấn công brute-force
\item Hệ thống được đồng bộ thời gian với NTP server
\item Chỉ các dịch vụ cần thiết cho hoạt động được kích hoạt
\item Giảm thiểu bề mặt tấn công (attack surface) của hệ thống
\end{itemize}

\subsection{Network Security}

\noindent \textbf{Configure network kernel parameters:}
\begin{verbatim}
sudo tee -a /etc/sysctl.d/99-cis.conf > /dev/null <<EOF
net.ipv4.ip_forward = 0
net.ipv4.conf.all.send_redirects = 0
net.ipv4.icmp_ignore_bogus_error_responses = 1
net.ipv4.icmp_echo_ignore_broadcasts = 1
net.ipv4.conf.all.accept_redirects = 0
net.ipv4.conf.all.secure_redirects = 0
net.ipv4.conf.all.rp_filter = 1
net.ipv4.conf.all.accept_source_route = 0
net.ipv4.conf.all.log_martians = 1
net.ipv4.tcp_syncookies = 1
net.ipv6.conf.all.accept_ra = 0
EOF

sudo sysctl -p /etc/sysctl.d/99-cis.conf                    # Apply all settings
\end{verbatim}

\noindent \textbf{Disable wireless services:}
\begin{verbatim}
sudo systemctl stop wpa_supplicant                     # Stop wireless service
sudo systemctl disable wpa_supplicant                  # Disable wireless boot
sudo systemctl stop NetworkManager                     # Stop NetworkManager
sudo systemctl disable NetworkManager                  # Disable NetworkManager boot
\end{verbatim}

\noindent \textbf{Configure IPv6 status:}
\begin{verbatim}
sudo tee -a /etc/sysctl.d/99-cis.conf > /dev/null <<EOF
net.ipv6.conf.all.disable_ipv6 = 1
net.ipv6.conf.default.disable_ipv6 = 1
EOF

sudo sysctl -p /etc/sysctl.d/99-cis.conf              # Apply IPv6 settings
\end{verbatim}

\noindent \textbf{Kết quả mong đợi:}
\begin{itemize}
\item Các tham số kernel mạng được cấu hình theo tiêu chuẩn bảo mật CIS
\item IP forwarding bị vô hiệu hóa để ngăn chặn router functionality
\item ICMP redirects và source routes bị chặn để chống lại spoofing attacks
\item IPv6 được vô hiệu hóa nếu không sử dụng
\item Các dịch vụ wireless không cần thiết được vô hiệu hóa
\item TCP SYN cookies được kích hoạt để chống lại SYN flood attacks
\end{itemize}

\subsection{Firewall Configuration}

\noindent \textbf{Configure UFW firewall:}
\begin{verbatim}
sudo apt install -y ufw                                    # Install UFW firewall
sudo apt remove --purge iptables-persistent                # Remove conflicting firewall
sudo ufw enable                                            # Enable firewall
sudo ufw default deny incoming                             # Default deny inbound
sudo ufw default allow outgoing                            # Default allow outbound
sudo ufw allow in on lo                                    # Allow loopback traffic
sudo ufw allow 22/tcp                                      # Allow SSH access
sudo ufw allow 80/tcp                                      # Allow HTTP (if needed)
sudo ufw allow 443/tcp                                     # Allow HTTPS (if needed)
sudo ufw status verbose                                    # Check firewall status
\end{verbatim}

\noindent \textbf{Configure firewall logging:}
\begin{verbatim}
sudo ufw logging on                                       # Enable firewall logging
sudo ufw logging medium                                   # Set log level to medium
sudo tail -f /var/log/ufw.log                             # Monitor firewall logs
\end{verbatim}

\noindent \textbf{Advanced firewall rules:}
\begin{verbatim}
sudo ufw deny 23/tcp                                      # Block telnet
sudo ufw deny 135/tcp                                     # Block RPC
sudo ufw deny 139/tcp                                     # Block NetBIOS
sudo ufw deny 445/tcp                                     # Block SMB
sudo ufw deny 1433/tcp                                    # Block SQL Server
sudo ufw deny 3389/tcp                                    # Block RDP
sudo ufw reload                                           # Apply all rules
\end{verbatim}

\noindent \textbf{Kết quả mong đợi:}
\begin{itemize}
\item UFW firewall được kích hoạt và bảo vệ hệ thống
\item Policy mặc định: deny inbound, allow outbound
\item Các cổng nguy hiểm (telnet, RPC, NetBIOS, etc.) bị chặn
\item Chỉ các cổng cần thiết (SSH, HTTP, HTTPS) được cho phép
\item Logging được kích hoạt để monitor và audit traffic
\item Loopback traffic được cho phép cho các dịch vụ nội bộ
\end{itemize}

\subsection{Access Control}

\noindent \textbf{SSH Security Hardening:}
\begin{verbatim}
sudo chmod 600 /etc/ssh/sshd_config                         # Secure SSH config
sudo chmod 600 /etc/ssh/ssh_host_*_key                      # Secure private keys
sudo chmod 644 /etc/ssh/ssh_host_*_key.pub                  # Secure public keys

sudo tee -a /etc/ssh/sshd_config > /dev/null <<EOF
PermitEmptyPasswords no
PermitRootLogin no
MaxAuthTries 3
MaxSessions 3
HostbasedAuthentication no
GSSAPIAuthentication no
ClientAliveInterval 300
ClientAliveCountMax 0
EOF

sudo systemctl restart sshd                                  # Apply SSH changes
\end{verbatim}

\noindent \textbf{Configure PAM Password Quality:}
\begin{verbatim}
sudo apt install -y libpam-pwquality                        # Install password quality

sudo tee -a /etc/pam.d/common-password > /dev/null <<EOF
password requisite pam_pwquality.so retry=3 \
minlen=8 
ucredit=-1 
lcredit=-1 
dcredit=-1 
ocredit=-1
EOF

sudo tee -a /etc/security/pwquality.conf > /dev/null <<EOF
minlen=8
ucredit=-1
lcredit=-1
dcredit=-1
ocredit=-1
EOF
\end{verbatim}

\noindent \textbf{Configure SUDO Security:}
\begin{verbatim}
sudo tee -a /etc/sudoers.d/security > /dev/null <<EOF
Defaults timestamp_timeout=15
Defaults passwd_tries=3
Defaults use_pty
Defaults log_input, log_output
EOF

sudo chmod 440 /etc/sudoers.d/security                        # Secure sudoers file
\end{verbatim}

\noindent \textbf{Configure User Defaults:}
\begin{verbatim}
sudo tee -a /etc/login.defs > /dev/null <<EOF
UMASK 077
PASS_MAX_DAYS 90
PASS_MIN_DAYS 1
PASS_WARN_AGE 7
LOGIN_RETRIES 5
LOGIN_TIMEOUT 60
EOF
\end{verbatim}

\noindent \textbf{Kết quả mong đợi:}
\begin{itemize}
\item SSH được hardening với các thiết lập bảo mật tối đa
\item Root login bị vô hiệu hóa, chỉ cho phép SSH key authentication
\item Password quality enforcement được áp dụng với các yêu cầu phức tạp
\item Sudo được cấu hình với logging và timeout hợp lý
\item User defaults được cấu hình với các chính sách bảo mật chặt chẽ
\item Các file SSH keys được bảo vệ với permissions phù hợp
\end{itemize}

\subsection{Logging \& Auditing}

\noindent \textbf{Configure auditd service:}
\begin{verbatim}
sudo apt install -y auditd audispd-plugins                # Install audit packages
sudo systemctl enable auditd                             # Enable auditd boot
sudo systemctl start auditd                              # Start auditd service
sudo tee /etc/default/grub.d/99-audit.cfg > /dev/null <<EOF
quiet splash audit=1 audit_backlog_limit=8192
EOF
sudo update-grub                                         # Update bootloader
\end{verbatim}

\noindent \textbf{Configure auditd settings:}
\begin{verbatim}
sudo tee -a /etc/audit/auditd.conf > /dev/null <<EOF
max_log_file = 8
max_log_file_action = suspend
space_left_action = single
admin_space_left_action = suspend
EOF

sudo systemctl restart auditd                            # Apply auditd changes
\end{verbatim}

\noindent \textbf{Configure journald and rsyslog:}
\begin{verbatim}
sudo systemctl enable systemd-journald                     # Enable journald service
sudo systemctl start systemd-journald                      # Start journald service
sudo apt install -y rsyslog                               # Install rsyslog package
sudo systemctl enable rsyslog                             # Enable rsyslog service
sudo systemctl start rsyslog                              # Start rsyslog service

sudo tee -a /etc/systemd/journald.conf > /dev/null <<EOF
ForwardToSyslog=no
Compress=yes
Storage=persistent
EOF

sudo tee -a /etc/rsyslog.d/50-default.conf > /dev/null <<EOF
\$FileCreateMode 0640
EOF

sudo chmod 0755 /var/log && \                           # Secure log directory
sudo chown root:root /var/log
sudo systemctl restart systemd-journald rsyslog           # Apply logging changes
\end{verbatim}

\noindent \textbf{Configure AIDE for file integrity:}
\begin{verbatim}
sudo apt install -y aide                                   # Install AIDE
sudo aide --init                                          # Initialize AIDE database
sudo mv /var/lib/aide/aide.db.new /var/lib/aide/aide.db    # Move initial database
sudo aide --check                                         # Run integrity check
sudo crontab - > /dev/null <<EOF
0 2 * * * /usr/bin/aide --check
EOF
\end{verbatim}

\noindent \textbf{Kết quả mong đợi:}
\begin{itemize}
\item Auditd được kích hoạt và cấu hình để monitor system events
\item Logging system được thiết lập với persistent storage và appropriate permissions
\item AIDE database được khởi tạo để monitor file integrity
\item Daily integrity checks được schedule để phát hiện thay đổi bất thường
\item Log files được bảo vệ với permissions phù hợp
\item System events được ghi lại để phục vụ cho forensic analysis
\end{itemize}

\subsection{System Maintenance}

\noindent \textbf{Configure critical system file permissions:}
\begin{verbatim}
sudo chmod 0644 /etc/passwd /etc/passwd- /etc/group /etc/group- /etc/shells
sudo chown root:root /etc/passwd /etc/passwd- /etc/group /etc/group- /etc/shells
sudo chmod 0000 /etc/shadow /etc/shadow- /etc/gshadow /etc/gshadow-
sudo chown root:shadow /etc/shadow /etc/shadow- /etc/gshadow /etc/gshadow-
sudo chmod 0600 /etc/security/opasswd
sudo chown root:root /etc/security/opasswd
\end{verbatim}

\noindent \textbf{Find and secure world writable files:}
\begin{verbatim}
find / -type f -perm -002 \                              # Find world writable files
     -not -path "/proc/*" \                              # Exclude /proc directory
     -not -path "/sys/*" \                               # Exclude /sys directory  
     -not -path "/dev/*" 2>/dev/null                     # Exclude /dev directory

find / -type d -perm -002 \                              # Find world writable dirs
     -not -path "/proc/*" \                              # Exclude /proc directory
     -not -path "/sys/*" \                               # Exclude /sys directory
     -not -path "/dev/*" 2>/dev/null                     # Exclude /dev directory
     
chmod o-w /etc/passwd /etc/shadow /etc/group \           # Remove world write
     /etc/gshadow 2>/dev/null || true
\end{verbatim}

\noindent \textbf{Review SUID/SGID files:}
\begin{verbatim}
find / -type f -perm -4000 \                             # Find SUID files
     -not -path "/proc/*" \                              # Exclude /proc directory
     -not -path "/sys/*" \                               # Exclude /sys directory
     -not -path "/dev/*" 2>/dev/null                     # Exclude /dev directory
     
find / -type f -perm -2000 \                             # Find SGID files
     -not -path "/proc/*" \                              # Exclude /proc directory
     -not -path "/sys/*" \                               # Exclude /sys directory
     -not -path "/dev/*" 2>/dev/null                     # Exclude /dev directory
     
chmod 0755 /usr/bin/passwd 2>/dev/null || true           # Secure passwd binary
chmod 0755 /usr/bin/gpasswd 2>/dev/null || true          # Secure gpasswd binary
chmod 0755 /usr/bin/chsh 2>/dev/null || true             # Secure chsh binary
\end{verbatim}

\noindent \textbf{Configure home directory permissions:}
\begin{verbatim}
find /home -type d -exec chmod 0750 {} \; \              # Set directory permissions
     2>/dev/null || true
find /home -type f -exec chmod 0640 {} \; \              # Set file permissions
     2>/dev/null || true
chmod 0755 /home 2>/dev/null || true                     # Set home base permission
\end{verbatim}

\noindent \textbf{System integrity checks:}
\begin{verbatim}
sudo debsums -c                                          # Check package integrity
sudo dpkg --audit                                        # Check for broken packages
sudo apt --fix-broken install                            # Fix broken packages
sudo apt autoremove                                      # Remove unused packages
sudo apt autoclean                                       # Clean package cache
\end{verbatim}

\noindent \textbf{Kết quả mong đợi:}
\begin{itemize}
\item Các file hệ thống quan trọng được bảo vệ với permissions phù hợp
\item World writable files và directories được xác định và secured
\item SUID/SGID files được review và chỉ giữ lại những file cần thiết
\item Home directories được cấu hình với permissions an toàn
\item Package integrity được kiểm tra và các broken packages được fix
\item System được dọn dẹp với unused packages và cache removal
\item File system permissions tuân thủ theo tiêu chuẩn bảo mật
\end{itemize}

\section{Automation}

\subsection{Initial Setup Automation}

\noindent \textbf{Optimized AppArmor Configuration:}
\begin{verbatim}
- name: Configure AppArmor security framework
  block:
    - name: Install AppArmor packages
      package:
        name: "{{ item }}"
        state: present
      loop:
        - apparmor
        - apparmor-utils

    - name: Enable AppArmor in bootloader
      lineinfile:
        path: /etc/default/grub
        regexp: '^GRUB_CMDLINE_LINUX_DEFAULT='
        line: 'GRUB_CMDLINE_LINUX_DEFAULT="quiet splash apparmor=1 security=apparmor"'

    - name: Start and enable AppArmor service
      service:
        name: apparmor
        state: started
        enabled: yes

    - name: Enforce all AppArmor profiles
      shell: aa-enforce /etc/apparmor.d/*
      register: apparmor_enforce
      changed_when: false
      failed_when: false
\end{verbatim}

\noindent \textbf{Optimized Filesystem Security:}
\begin{verbatim}
- name: Configure filesystem security
  block:
    - name: Secure temporary directories
      shell: |
        mount -o remount,nodev,nosuid,noexec {{ item }} \
        2>/dev/null || true
      loop:
        - /tmp
        - /dev/shm
      register: tmp_secured
      changed_when: tmp_secured.rc == 0

    - name: Add /var/tmp to fstab with security options
      lineinfile:
        path: /etc/fstab
        regexp: '^.*/var/tmp.*'
        line: "tmpfs /var/tmp tmpfs defaults,nodev,nosuid,noexec 0 0"
        create: yes

    - name: Mount /var/tmp with security options
      command: mount -t tmpfs -o nodev,nosuid,noexec tmpfs /var/tmp
      register: var_tmp_mount
      failed_when: var_tmp_mount.rc != 0 and "already mounted" not in var_tmp_mount.stderr
\end{verbatim}

\noindent \textbf{Optimized Login Banners:}
\begin{verbatim}
- name: Configure login banners and SSH
  block:
    - name: Create login warning banners
      copy:
        content: "Authorized uses only. All activity may be monitored and reported.\n"
        dest: "{{ item }}"
        owner: root
        group: root
        mode: '0644'
      loop:
        - /etc/issue
        - /etc/issue.net
        - /etc/motd

    - name: Configure SSH to display banner
      lineinfile:
        path: /etc/ssh/sshd_config
        regexp: '^Banner'
        line: 'Banner /etc/issue.net'
        state: present
      notify: restart sshd
\end{verbatim}

\noindent \textbf{Kết quả mong đợi:}
\begin{itemize}
\item AppArmor được tự động cài đặt và cấu hình với enforce mode
\item Filesystem security được áp dụng một cách nhất quán trên tất cả systems
\item Login banners được triển khai đồng bộ với nội dung thống nhất
\item Quá trình automation đảm bảo consistency và repeatability
\item Các tasks được thực hiện với error handling và appropriate logging
\item Configuration drift được ngăn chặn thông qua standardized automation
\end{itemize}

\subsection{Services Automation}

\noindent \textbf{Optimized Service Management:}
\begin{verbatim}
- name: Disable unnecessary services
  systemd:
    name: "{{ item }}"
    state: stopped
    enabled: no
  loop:
    - telnet.socket
    - telnet
    - rsh.socket
    - rsh
    - rlogin.socket
    - rlogin
    - ypbind
    - tftp.socket
    - tftp
    - isc-dhcp-server
    - slapd
    - apache2
    - nfs-server
    - rpcbind
    - cups
    - avahi-daemon
    - bluetooth
  when: ansible_facts.services[item] is defined
  ignore_errors: yes
\end{verbatim}

\noindent \textbf{Optimized Fail2ban Configuration:}
\begin{verbatim}
- name: Configure fail2ban intrusion prevention
  block:
    - name: Install and enable fail2ban
      apt:
        name: fail2ban
        state: present
      notify: restart fail2ban

    - name: Configure fail2ban settings
      copy:
        content: |
          [DEFAULT]
          bantime = 3600
          findtime = 600
          maxretry = 3
          
          [sshd]
          enabled = true
          port = ssh
          logpath = /var/log/auth.log
          maxretry = 3
        dest: /etc/fail2ban/jail.local
        owner: root
        group: root
        mode: '0644'
      notify: restart fail2ban
\end{verbatim}

\noindent \textbf{Optimized Time Synchronization:}
\begin{verbatim}
- name: Configure time synchronization
  block:
    - name: Install NTP service
      apt:
        name: ntp
        state: present

    - name: Enable and start NTP
      service:
        name: ntp
        state: started
        enabled: yes

    - name: Enable system time synchronization
      command: timedatectl set-ntp true
      changed_when: false
\end{verbatim}

\noindent \textbf{Kết quả mong đợi:}
\begin{itemize}
\item Các dịch vụ không cần thiết được tự động vô hiệu hóa trên toàn bộ infrastructure
\item Fail2ban được triển khai với configuration chuẩn hóa
\item Time synchronization được đảm bảo trên tất cả systems
\item Service management được thực hiện với proper error handling
\item Automation đảm bảo consistency across multiple environments
\item Các services được monitor và maintained trong desired state
\end{itemize}

\subsection{Network Security Automation}

\noindent \textbf{Optimized Network Parameters Configuration:}
\begin{verbatim}
- name: Configure network security parameters
  sysctl:
    name: "{{ item.name }}"
    value: "{{ item.value }}"
    state: present
    sysctl_file: /etc/sysctl.d/99-cis.conf
  loop:
    - { name: net.ipv4.ip_forward, value: '0' }
    - { name: net.ipv4.conf.all.send_redirects, value: '0' }
    - { name: net.ipv4.icmp_ignore_bogus_error_responses, value: '1' }
    - { name: net.ipv4.icmp_echo_ignore_broadcasts, value: '1' }
    - { name: net.ipv4.conf.all.accept_redirects, value: '0' }
    - { name: net.ipv4.conf.all.secure_redirects, value: '0' }
    - { name: net.ipv4.conf.all.rp_filter, value: '1' }
    - { name: net.ipv4.conf.all.accept_source_route, value: '0' }
    - { name: net.ipv4.conf.all.log_martians, value: '1' }
    - { name: net.ipv4.tcp_syncookies, value: '1' }
    - { name: net.ipv6.conf.all.accept_ra, value: '0' }

- name: Apply network settings
  command: sysctl -p /etc/sysctl.d/99-cis.conf
  changed_when: false
\end{verbatim}

\noindent \textbf{Optimized Wireless Services Management:}
\begin{verbatim}
- name: Disable wireless services
  systemd:
    name: "{{ item }}"
    state: stopped
    enabled: no
  loop:
    - wpa_supplicant
    - NetworkManager
  when: ansible_facts.services[item] is defined
  ignore_errors: yes
\end{verbatim}

\noindent \textbf{Optimized IPv6 Configuration:}
\begin{verbatim}
- name: Configure IPv6 settings
  sysctl:
    name: "{{ item }}"
    value: '1'
    state: present
    sysctl_file: /etc/sysctl.d/99-cis.conf
  loop:
    - net.ipv6.conf.all.disable_ipv6
    - net.ipv6.conf.default.disable_ipv6
\end{verbatim}

\noindent \textbf{Kết quả mong đợi:}
\begin{itemize}
\item Network security parameters được áp dụng một cách nhất quán
\item IPv6 được quản lý centralized thông qua automation
\item Wireless services được vô hiệu hóa tự động khi không cần thiết
\item Sysctl configuration được version controlled và auditable
\item Network hardening được thực hiện với minimal manual intervention
\item Changes được tracked và rollback capability được đảm bảo
\end{itemize}

\subsection{Firewall Automation}

\noindent \textbf{Optimized UFW Firewall Configuration:}
\begin{verbatim}
- name: Configure UFW firewall
  block:
    - name: Install UFW and remove conflicting packages
      package:
        name: "{{ item }}"
        state: "{{ item.state | default('present') }}"
      loop:
        - { name: ufw, state: present }
        - { name: iptables-persistent, state: absent }

    - name: Enable and start UFW service
      service:
        name: ufw
        state: started
        enabled: yes

    - name: Configure UFW default policies
      ufw:
        state: "{{ item.state | default('enabled') }}"
        policy: "{{ item.policy }}"
        direction: "{{ item.direction }}"
      loop:
        - { state: enabled, policy: deny, direction: incoming }
        - { state: enabled, policy: allow, direction: outgoing }

    - name: Configure UFW loopback traffic
      ufw:
        rule: allow
        interface: lo
        direction: in

    - name: Configure essential UFW rules
      ufw:
        rule: "{{ item.rule }}"
        port: "{{ item.port }}"
        proto: "{{ item.proto | default('tcp') }}"
      loop:
        - { rule: allow, port: 22, proto: tcp }      # SSH
        - { rule: deny, port: 23, proto: tcp }       # Block telnet
        - { rule: deny, port: 135, proto: tcp }      # Block RPC
        - { rule: deny, port: 139, proto: tcp }      # Block NetBIOS
        - { rule: deny, port: 445, proto: tcp }      # Block SMB
        - { rule: deny, port: 3389, proto: tcp }     # Block RDP
\end{verbatim}

\noindent \textbf{Optimized Firewall Logging Configuration:}
\begin{verbatim}
- name: Configure firewall logging
  block:
    - name: Enable UFW logging
      ufw:
        logging: on

    - name: Set UFW log level
      lineinfile:
        path: /etc/ufw/ufw.conf
        regexp: '^LOGLEVEL='
        line: 'LOGLEVEL=medium'
      notify: restart ufw
\end{verbatim}

\noindent \textbf{Kết quả mong đợi:}
\begin{itemize}
\item UFW firewall được tự động cấu hình với standardized rules
\item Firewall policies được áp dụng consistent across environments
\item Logging được kích hoạt với appropriate log levels
\item Rules được managed through code (Infrastructure as Code)
\item Changes được tracked và audit trail được maintained
\item Firewall configuration được version controlled và testable
\end{itemize}

\subsection{Access Control Automation}

\noindent \textbf{Optimized SSH Security Configuration:}
\begin{verbatim}
- name: Configure SSH security
  block:
    - name: Secure SSH configuration files
      file:
        path: "{{ item.path }}"
        mode: "{{ item.mode }}"
        owner: root
        group: root
      loop:
        - { path: /etc/ssh/sshd_config, mode: '0600' }
        - { path: /etc/ssh/ssh_host_rsa_key, mode: '0600' }
        - { path: /etc/ssh/ssh_host_ecdsa_key, mode: '0600' }
        - { path: /etc/ssh/ssh_host_ed25519_key, mode: '0600' }
        - { path: /etc/ssh/ssh_host_rsa_key.pub, mode: '0644' }
        - { path: /etc/ssh/ssh_host_ecdsa_key.pub, mode: '0644' }
        - { path: /etc/ssh/ssh_host_ed25519_key.pub, mode: '0644' }

    - name: Configure SSH security parameters
      lineinfile:
        path: /etc/ssh/sshd_config
        regexp: '^#?{{ item.key }}'
        line: '{{ item.key }} {{ item.value }}'
        state: present
      loop:
        - { key: PermitEmptyPasswords, value: no }
        - { key: PermitRootLogin, value: no }
        - { key: MaxAuthTries, value: 3 }
        - { key: MaxSessions, value: 3 }
        - { key: HostbasedAuthentication, value: no }
        - { key: GSSAPIAuthentication, value: no }
        - { key: ClientAliveInterval, value: 300 }
        - { key: ClientAliveCountMax, value: 0 }
      notify: restart sshd
\end{verbatim}

\noindent \textbf{Optimized PAM Password Quality Configuration:}
\begin{verbatim}
- name: Configure PAM password quality
  block:
    - name: Install PAM password quality module
      apt:
        name: libpam-pwquality
        state: present

    - name: Configure password policy in PAM
      lineinfile:
        path: /etc/pam.d/common-password
        line: "password requisite pam_pwquality.so retry=3 minlen=8 ucredit=-1 lcredit=-1 dcredit=-1 ocredit=-1"

    - name: Configure password quality settings
      lineinfile:
        path: /etc/security/pwquality.conf
        line: "{{ item }}"
      loop:
        - "minlen=8"
        - "ucredit=-1"
        - "lcredit=-1"
        - "dcredit=-1"
        - "ocredit=-1"
\end{verbatim}

\noindent \textbf{Optimized SUDO Security Configuration:}
\begin{verbatim}
- name: Configure SUDO security
  block:
    - name: Configure SUDO security defaults
      lineinfile:
        path: /etc/sudoers.d/security
        line: "{{ item }}"
        create: yes
        mode: '0440'
      loop:
        - "Defaults timestamp_timeout=15"
        - "Defaults passwd_tries=3"
        - "Defaults use_pty"
        - "Defaults log_input, log_output"
\end{verbatim}

\noindent \textbf{Optimized User Defaults Configuration:}
\begin{verbatim}
- name: Configure user security defaults
  lineinfile:
    path: /etc/login.defs
    line: "{{ item }}"
  loop:
    - "UMASK 077"
    - "PASS_MAX_DAYS 90"
    - "PASS_MIN_DAYS 1"
    - "PASS_WARN_AGE 7"
    - "LOGIN_RETRIES 5"
    - "LOGIN_TIMEOUT 60"
\end{verbatim}

\noindent \textbf{Kết quả mong đợi:}
\begin{itemize}
\item SSH security được tự động hardening với consistent configuration
\item Password policies được áp dụng standardized across all users
\item Sudo configuration được managed centrally với proper logging
\item User defaults được cấu hình với security-first approach
\item Access control policies được enforced through automation
\item Configuration drift được eliminated và compliance được đảm bảo
\end{itemize}

\subsection{Logging \& Auditing Automation}

\noindent \textbf{Optimized Auditd Configuration:}
\begin{verbatim}
- name: Configure auditd service
  block:
    - name: Install audit packages
      apt:
        name:
          - auditd
          - audispd-plugins
        state: present

    - name: Enable and start auditd
      service:
        name: auditd
        state: started
        enabled: yes

    - name: Configure kernel auditing parameters
      lineinfile:
        path: /etc/default/grub
        regexp: '^GRUB_CMDLINE_LINUX_DEFAULT='
        line: 'GRUB_CMDLINE_LINUX_DEFAULT="quiet splash audit=1 audit_backlog_limit=8192"'

    - name: Configure auditd settings
      lineinfile:
        path: /etc/audit/auditd.conf
        regexp: '^{{ item.key }}'
        line: '{{ item.key }} = {{ item.value }}'
      loop:
        - { key: max_log_file, value: 8 }
        - { key: max_log_file_action, value: suspend }
        - { key: space_left_action, value: single }
        - { key: admin_space_left_action, value: suspend }
      notify: restart auditd
\end{verbatim}

\noindent \textbf{Optimized Logging System Configuration:}
\begin{verbatim}
- name: Configure logging system
  block:
    - name: Install and enable rsyslog
      apt:
        name: rsyslog
        state: present

    - name: Enable and start logging services
      service:
        name: "{{ item }}"
        state: started
        enabled: yes
      loop:
        - systemd-journald
        - rsyslog

    - name: Configure journald settings
      lineinfile:
        path: /etc/systemd/journald.conf
        regexp: '^#?{{ item.key }}'
        line: '{{ item.key }}={{ item.value }}'
      loop:
        - { key: ForwardToSyslog, value: no }
        - { key: Compress, value: yes }
        - { key: Storage, value: persistent }

    - name: Configure rsyslog file permissions
      lineinfile:
        path: /etc/rsyslog.d/50-default.conf
        regexp: '^\$FileCreateMode'
        line: '$FileCreateMode 0640'

    - name: Secure log directory
      file:
        path: /var/log
        mode: '0755'
        owner: root
        group: root
        recurse: yes
\end{verbatim}

\noindent \textbf{Optimized AIDE Configuration:}
\begin{verbatim}
- name: Configure AIDE file integrity monitoring
  block:
    - name: Install AIDE
      apt:
        name: aide
        state: present

    - name: Initialize AIDE database
      command: aide --init
      register: aide_init
      changed_when: false

    - name: Move AIDE database to production location
      command: mv /var/lib/aide/aide.db.new /var/lib/aide/aide.db
      when: aide_init.rc == 0

    - name: Schedule daily AIDE checks
      cron:
        name: "AIDE integrity check"
        minute: "0"
        hour: "2"
        job: "/usr/bin/aide --check"
        user: root
\end{verbatim}

\noindent \textbf{Kết quả mong đợi:}
\begin{itemize}
\item Auditd được tự động cấu hình với comprehensive logging
\item Logging system được standardized với persistent storage
\item AIDE database được initialized và scheduled checks được automated
\item Log retention policies được applied consistently
\item Audit trail được maintained cho compliance và forensic purposes
\item Monitoring và alerting được integrated với logging infrastructure
\end{itemize}

\subsection{System Maintenance Automation}

\noindent \textbf{Optimized System File Permissions:}
\begin{verbatim}
- name: Configure critical system file permissions
  file:
    path: "{{ item.path }}"
    mode: "{{ item.mode }}"
    owner: "{{ item.owner }}"
    group: "{{ item.group }}"
  loop:
    - { path: /etc/passwd, mode: '0644', owner: root, group: root }
    - { path: /etc/passwd-, mode: '0644', owner: root, group: root }
    - { path: /etc/group, mode: '0644', owner: root, group: root }
    - { path: /etc/group-, mode: '0644', owner: root, group: root }
    - { path: /etc/shadow, mode: '0000', owner: root, group: shadow }
    - { path: /etc/shadow-, mode: '0000', owner: root, group: shadow }
    - { path: /etc/gshadow, mode: '0000', owner: root, group: shadow }
    - { path: /etc/gshadow-, mode: '0000', owner: root, group: shadow }
    - { path: /etc/shells, mode: '0644', owner: root, group: root }
    - { path: /etc/security/opasswd, mode: '0600', owner: root, group: root }
\end{verbatim}

\noindent \textbf{Optimized World Writable Files Security:}
\begin{verbatim}
- name: Secure world writable files and directories
  block:
    - name: Find world writable files
      shell: |
        find / -type f -perm -002 \
        -not -path "/proc/*" \
        -not -path "/sys/*" \
        -not -path "/dev/*" \
        2>/dev/null || true
      register: world_writable_files
      changed_when: false

    - name: Display world writable files
      debug:
        msg: "World writable files: {{ world_writable_files.stdout_lines }}"
      when: world_writable_files.stdout_lines is defined

    - name: Find world writable directories
      shell: |
        find / -type d -perm -002 \
        -not -path "/proc/*" \
        -not -path "/sys/*" \
        -not -path "/dev/*" \
        2>/dev/null || true
      register: world_writable_dirs
      changed_when: false

    - name: Display world writable directories
      debug:
        msg: "World writable directories: {{ world_writable_dirs.stdout_lines }}"
      when: world_writable_dirs.stdout_lines is defined

    - name: Remove world writable permissions from critical files
      shell: |
        for file in /etc/passwd /etc/shadow /etc/group /etc/gshadow; do
          if [ -f "$file" ]; then
            chmod o-w "$file" 2>/dev/null || true
          fi
        done
      changed_when: false
\end{verbatim}

\noindent \textbf{Optimized SUID/SGID Review:}
\begin{verbatim}
- name: Review and secure SUID/SGID files
  block:
    - name: Find SUID files
      shell: |
        find / -type f -perm -4000 \
        -not -path "/proc/*" \
        -not -path "/sys/*" \
        -not -path "/dev/*" \
        2>/dev/null || true
      register: suid_files
      changed_when: false

    - name: Find SGID files
      shell: |
        find / -type f -perm -2000 \
        -not -path "/proc/*" \
        -not -path "/sys/*" \
        -not -path "/dev/*" \
        2>/dev/null || true
      register: sgid_files
      changed_when: false

    - name: Secure common SUID binaries
      file:
        path: "{{ item }}"
        mode: '0755'
      loop:
        - /usr/bin/passwd
        - /usr/bin/gpasswd
        - /usr/bin/chsh
      ignore_errors: yes
\end{verbatim}

\noindent \textbf{Optimized System Integrity Maintenance:}
\begin{verbatim}
- name: Perform system integrity checks
  block:
    - name: Check package integrity
      command: debsums -c
      register: debsums_check
      changed_when: false
      ignore_errors: yes

    - name: Audit broken packages
      command: dpkg --audit
      register: dpkg_audit
      changed_when: false
      ignore_errors: yes

    - name: Fix broken packages
      apt:
        upgrade: dist
        fix_broken: yes
      when: dpkg_audit.rc != 0

    - name: Remove unused packages
      apt:
        autoremove: yes
        purge: yes

    - name: Clean package cache
      apt:
        autoclean: yes
\end{verbatim}

\noindent \textbf{Kết quả mong đợi:}
\begin{itemize}
\item System file permissions được tự động cấu hình và maintained
\item World writable files được identified và remediated automatically
\item SUID/SGID files được reviewed và secured với minimal manual intervention
\item System integrity checks được performed regularly
\item Package management được automated với proper cleanup
\item Configuration drift được prevented và compliance được ensured
\item Maintenance tasks được scheduled và executed consistently
\end{itemize}

